% ---------------------------------------------------------------
%  Frederick Soddy, CARTESIAN ECONOMICS (Hendersons, London, 1922)
%  A reset of the 1922 edition from the Harvard/Widener scan
%  (barcode 32044010374809).  Pagination follows the original.
%  Set in MFB Oldstyle = Century Oldstyle (M. F. Benton, ATF 1909).
% ---------------------------------------------------------------
\documentclass[10.5pt,twoside]{book}

\usepackage{fontspec}
\usepackage[
  paperwidth=6in, paperheight=9in,
  inner=1.00in, outer=0.82in,
  top=0.72in, bottom=0.80in,
  headsep=15pt, footskip=22pt
]{geometry}
\usepackage{microtype}
\usepackage{lettrine}
\usepackage{fancyhdr}
\usepackage[hidelinks,pdfusetitle]{hyperref}
\hypersetup{
  pdftitle={Cartesian Economics: The Bearing of Physical Science upon State Stewardship},
  pdfauthor={Frederick Soddy},
  pdfsubject={Two lectures, Birkbeck College and the London School of Economics, November 1921},
  pdfkeywords={Soddy, thermodynamics, economics, wealth, debt, energy, 1922, Hendersons}}
\usepackage{ragged2e}
\usepackage[normalem]{ulem}

\setmainfont{MFBOldstyle}[
  Extension      = .otf,
  UprightFont    = *-Regular,
  ItalicFont     = *-Italic,
  BoldFont       = *-Bold,
  Ligatures      = TeX,
  Numbers        = OldStyle
]
\newfontfamily\liningfont{MFBOldstyle}[
  Extension=.otf, UprightFont=*-Regular, ItalicFont=*-Italic,
  BoldFont=*-Bold, Ligatures=TeX, Numbers=Lining]

\linespread{1.06}
\setlength{\parindent}{1.4em}
\raggedbottom
\frenchspacing
\hyphenpenalty=150
\tolerance=2500
\emergencystretch=1.5em
\clubpenalty=3000
\widowpenalty=3000

% --- letterspaced caps, the way a 1922 comp would have done it ----
\newcommand{\lsp}[2]{{\addfontfeature{LetterSpace=#1}\MakeUppercase{#2}}}
\newcommand{\scaps}[1]{{\addfontfeature{LetterSpace=4}\textsc{#1}}}

% --- running heads: folio outside, title centred --------------------
\pagestyle{fancy}
\fancyhf{}
\renewcommand{\headrulewidth}{0pt}
\fancyhead[LE]{\normalsize\thepage\hfill}
\fancyhead[RO]{\hfill\normalsize\thepage}
\fancyhead[C]{{\addfontfeature{LetterSpace=8}\normalsize CARTESIAN ECONOMICS.}}
\fancypagestyle{plain}{\fancyhf{}\renewcommand{\headrulewidth}{0pt}}

% --- block quotations: smaller, indented both sides -----------------
\newenvironment{extract}
  {\par\addvspace{7pt}\begingroup\small\linespread{1.02}\selectfont
   \leftskip=1.9em \rightskip=1.2em \parindent=1.2em \noindent\ignorespaces}
  {\par\endgroup\addvspace{7pt}}

% --- forced original page breaks -----------------------------------
\newcommand{\origpage}[1]{\clearpage\setcounter{page}{#1}}

% --- footnote styling ----------------------------------------------
\usepackage[flushmargin]{footmisc}
\renewcommand{\thefootnote}{*}
\renewcommand{\footnotelayout}{\footnotesize}
\setlength{\skip\footins}{14pt plus 3pt}
\renewcommand{\footnoterule}{\vspace*{-3pt}\rule{6em}{0.4pt}\vspace*{2.5pt}}

\begin{document}
\frontmatter
\pagestyle{empty}

% ================= HALF TITLE =================
\vspace*{2.1in}
\begin{center}
  {\addfontfeature{LetterSpace=14}\large CARTESIAN ECONOMICS}
\end{center}
\clearpage\null\clearpage

% ================= TITLE PAGE =================
\vspace*{1.45in}
\begin{center}
  {\addfontfeature{LetterSpace=10}\fontsize{21}{25}\selectfont Cartesian Economics}\\[26pt]
  {\fontsize{12.5}{17}\selectfont\itshape The Bearing of Physical Science\\
   upon State Stewardship.\ \ By\\[3pt]
   Frederick Soddy, M.A., F.R.S.}\\[16pt]
  \rule{1.5in}{0.5pt}\\[2.30in]
  {\addfontfeature{LetterSpace=16}\footnotesize LONDON:}\\[7pt]
  {\addfontfeature{LetterSpace=14}\fontsize{15}{18}\selectfont HENDERSONS}\\[5pt]
  {\addfontfeature{LetterSpace=10}\footnotesize 66 CHARING CROSS ROAD}\\[12pt]
\end{center}
\clearpage

% ================= IMPRINT =================
\null\vfill
\begin{center}
  {\footnotesize First published, May, 1922}\\[4pt]
  {\footnotesize Reprinted, October, 1922}
\end{center}
\vspace*{1.4in}
\clearpage

% ================= TEXT =================
\mainmatter
\pagestyle{fancy}
\setcounter{page}{3}
\thispagestyle{plain}

\begin{center}
  {\addfontfeature{LetterSpace=8}\fontsize{16}{20}\selectfont Cartesian Economics.}\\[13pt]
  {\addfontfeature{LetterSpace=6}\footnotesize THE BEARING OF PHYSICAL SCIENCE UPON STATE\\[2pt]
   STEWARDSHIP.}\\[13pt]
  {\footnotesize\itshape Two Lectures to the Student Unions of Birkbeck College and\\
   The London School of Economics, November 10th and 17th,\\
   1921.\footnotemark}\\[9pt]
  {\footnotesize By Frederick Soddy, M.A., F.R.S.}\\[16pt]
  {\addfontfeature{LetterSpace=6}\footnotesize FIRST LECTURE.\,---\,Chairman, Sir Richard Gregory.}\\[10pt]
\end{center}
\footnotetext{These lectures owe much to a long controversial correspondence
carried on with Dr.\ H. Lyster Jameson, whose sad death is just reported in
the papers. Dr.\ Jameson upheld the neo-Marxian or proletarian view in
economics and the determinist or, as I style it, the ``ultra-materialistic,''
philosophy, and from the controversy my own views gained in definiteness and
clearness.\par\vspace{2pt}\hfill\textit{March 3rd, 1922.}\quad F. S.}

\lettrine[lines=2,findent=2pt,nindent=0pt,slope=0pt]{I}{T} is my intention to
try to bring the existing knowledge of the physical sciences to bear upon the
question ``How do men live?'' This question ought to be the first the
economist should try to answer. I am by no means the first to essay this task,
but the modern economist seems to have forgotten that there is such a
question, whilst the earlier ones lived at a stage of the development of
scientific knowledge when no exact answer was forthcoming.

My own point of departure could not be better illustrated than by a quotation
from Descartes, and the aspects I propose to examine might well be called
``Cartesian Economics.''

\begin{extract}
``Starting from the forms of knowledge most useful to life, instead of from
that speculative philosophy taught in our schools, and knowing the force and
processes of fire, the air, the stars and all the other bodies which surround
us as distinctly as we know the different occupations of our own workmen, we
shall be able to employ them in the same fashion and so render ourselves as
the masters and possessors of nature and contribute to the perfection of the
human life.''
\end{extract}

\origpage{4}

The enormous progress made in the mastery of man over nature and the meagre
contribution to the perfection of the human life is a contrast that can only
be accounted for by some such enquiry as that which I propose to essay. But in
language more homely than that of Descartes I may illustrate my starting point
by means of a story. An expert organist, drawing enthusiastic applause from
his audience, was surprised and annoyed by the blower coming to the front of
the screen and remarking to him, ``Yes! we played that piece very well.'' The
blower not being encouraged in well-doing, in the next piece the divine music
rose majestically to its climax and then petered out in a dismal wail, whilst
a head appeared round the screen and remarked, ``Now! is it we?'' Nor is it
without significance to note that, since the occurrence, the human labour upon
which the organist relied has been replaced so completely by electric power.
Power, rather than any qualifying adjective, human, mechanical or electrical,
is the starting point of Cartesian economics.

At the risk of being redundant, let me illustrate what I mean by the question,
``How do men live?'' by asking what makes a railway train go. In one sense or
another the credit for the achievement may be claimed by the so-called
``engine-driver,'' the guard, the signalman, the manager, the capitalist, or
the share-holder,---or, again, by the scientific pioneers who discovered the
nature of fire, by the inventors who harnessed it, by Labour which built the
railway and the train. The fact remains that all of them by their united
efforts could not drive the train. The real engine-driver is the coal. So, in
the present state of science, the answer to the question how men live, or how
anything lives, or how inanimate nature lives, in the sense in which we speak
of the life of a waterfall or of any other manifestation of continued
liveliness, is, with few and unimportant exceptions, ``By sunshine.'' Switch
off the sun and a world would result lifeless, not only in the sense of
animate life, but also in respect of by far the greater part of the life of
inanimate nature. The volcanoes, as now, might occasionally erupt, the tides
would ebb and flow on an otherwise stagnant ocean, and the newly discovered
phenomena of radioactivity would persist. But it is sunshine which provides
the power not only of the winds and waters but also of every form of life yet
known. The starting point of Cartesian economics is thus the well-known laws
of the conservation and transformation of

\origpage{5}

\noindent energy, usually referred to as the first and second laws of
thermodynamics.

But let us, before starting on this quest, go to the other extreme and try to
obtain a consistent mental picture of the whole of knowledge and the
inter-relationship of the sciences, if only to rebut two, in my opinion,
errors, caricatured rather than described by the terms ``Mechanistic'' and
``Vitalistic'' in philosophy. In any classification of the sciences it is
customary to distinguish three great groups, (1) the Mechanical,
Physico-Chemical and Mathematical; (2) the Biological; and (3) the Mental,
inter-related much as in the order enumerated, like the links of a chain, two
end links and a middle link, and the latter, the great division of animal and
vegetable life, alone in direct relation to the two ends.

Without perhaps any more quantitative notion than that the subjects enumerated
in the list are in order of cognateness, one may regard the first link of the
chain as commencing with the present ultimate realities of
physics---Electricity, Energy, Ether, and Matter in increasing complexity from
the element to the complex colloid, actuated by the liveliest Brownian
movement under the microscope, and yet not alive. The second link begins with
the simplest unicellular form of life and stretches in ever-increasing
complexity from amoeba to man. The third begins with rudimentary forms of
instinctive behaviour closely allied to those produced by purely physical
stimuli in simple creatures, through free-will and the deliberate choice of
action to secure predetermined ends, up to the human reason and the highest
intellectual, aesthetic, ethical, moral and spiritual perceptions of humanity.

My own philosophy can be expressed by a line of Kipling---``East is East and
West is West and never the twain shall meet.'' One thousand years hence it is
certain, if civilisation lasts as long, that men will still be examining into
the fundamental realities of the physical world, far beyond the world of atoms
and energy which marks its present boundary. Equally, we may believe that they
may in psychological enquiry and in spiritual perceptions be much advanced
from our time. In each direction possibilities of further knowledge extend
\textit{ad infinitum}, but in each direction diametrically away from and not
towards the problems of life. It is in this middle field that economics lies,
unaffected whether by the ultimate philosophy of the

\origpage{6}

\noindent electron or of the soul, and concerned rather with the interaction
with the middle world of life of these two end worlds of physics and mind in
their commonest everyday aspects, matter and energy on the one hand, obeying
the laws of mathematical probability or chance as exhibited in the inanimate
universe, and, on the other, with the guidance, direction and willing of these
blind forces and processes to predetermined ends. The physicist claims that
his world of matter and energy exists as a reality independent of life, and
points to the laws of conservation to show that it is eternal, without
beginning and without end, and to the record of the rocks to show that it is
life, rather than the universe of nature, which began. The theologian and
religious philosopher claims under the name of Deity the independent and
eternal existence of the qualities of guidance, will and direction outside of
life, and points to this to account for the ascending scale of evolution and
the appearance of perceptions above the level of animal. I have no claim or
call to express an opinion on the reality of the existence of intelligence
apart from and outside of life. But that life is the expression of the
interaction of two totally distinct things represented by probability and
free-will is to me self-evident, though the ultimate nature of those two
different things will probably remain, a thousand years hence, as far off as
ever.

It is simple now to indicate what to my mind are the two errors that hinder
progress. Both are monistic obsessions due to the mind in its innate desire to
reduce everything to its simplest terms ending by trying to reduce everything
to its simplest term. The first links up the two ends of the chain running in
diametrically opposite directions into a grand circle, and so gets the
sublimated conceptions of the mental world inextricably mixed up with the
physical. The oriental philosophies and religions seem to have been freer from
the cruder forms of this confusion than our own. In the early forms deities
were given physical powers analogous to those of trinitrotoluene, as for
example the hammer of Thor and the thunderbolt of Jove. The idea that the
physical universe must, like life, have had a beginning and therefore a
creator still survives. Heaven to the ordinary man is at once the abode of
disembodied souls and of constellations which perform their evolutions with
such mathematical accuracy that events therein can be accurately predicted in
advance. Its most recent phase is the theosophistical enquiry by occult powers
into the internal

\origpage{7}

\noindent structure of the atom and the prevalent belief that the discovery of
wireless telegraphy lends strong support to the reality of telepathy.

The second error is perhaps more common in the sphere of economics. It may be
called ``Ultra-Materialism'' and is the attempt to derive the whole of the
phenomena of life by continuous evolution from the inanimate world. We begin
with a nebula of primordial material condensing into ever more complex forms,
first to the light and then the heavy elements, then to chemical compounds up
to the complex colloid. By a continuation of the same processes such a complex
results that it is continually decomposing and as continually regenerating
itself. The inanimate molecules begin to live and life then runs through its
course of evolution up to man. This may satisfy a biologist, but it fails to
satisfy me as a chemist. I cannot conceive of inanimate mechanism, obeying the
laws of probability, by any continued series of successive steps developing
the powers of choice and reproduction any more than I can envisage any
increase in the complexity of an engine resulting in the production of the
``engine-driver'' and the power of its reproducing itself. I shall be told
that this is a pontifical expression of personal opinion. Unfortunately,
however, for this argument, inanimate mechanism happens to be my special study
rather than that of the biologist. It is the invariable characteristic of all
shallow and pretentious philosophy to seek the explanation of insoluble
problems in some other field than that of which the philosopher has first-hand
acquaintance. The biologist has first-hand knowledge of animate mechanism and
seeks the origin of it in colloid chemistry. The test of the hypothesis is not
so much what the biologist as what the chemist has to say about it. The
difference to my mind between dead and living matter is much that between
Niagara Falls thirty years ago and now, and is not to be explained by the laws
which Niagara formerly obeyed, by the laws of pure probability, but by their
opposite, the operations of intelligence, as typified in their most
rudimentary form by Clerk-Maxwell's conception of the ``sorting-demon.''

Life, or animate mechanism, is essentially to my mind a dualism, and any
attempt to subordinate either partner is fatal. But the economist is
peculiarly liable to mistake for laws of nature the laws of human nature and
to dignify this complex of

\origpage{8}

\noindent thermodynamical and social phenomena with the term ``inexorable
economic law.''

Is it any wonder that such crude confusions, such triumphs of mental instincts
over reason, experience and common sense, have produced a general sterility of
constructive thought? I cannot do better to illustrate this than to quote
Stephen Leacock in his most serious mood, and if it be objected that he is a
humorist, I can only retort that he is a professor of economics. It is rather
the avowal of the combination that is uncommon.

\begin{extract}
``Our studies consist only in the long-drawn proof of the futility of our
search after knowledge effected by exposing the errors of the past. Philosophy
is the science which proves that we can know nothing of the soul. Medicine is
the science which tells us that we know nothing of the body. Political Economy
is that which teaches that we know nothing of the laws of Wealth, and Theology
is the critical history of those errors from which we deduce our ignorance of
God.

``When I sit and warm my hands, as best I may, at the little heap of embers
that is now Political Economy, I cannot but contrast its dying glow with the
vainglorious and triumphant science that once it was.''
\end{extract}

Against this I would put the paradoxical words of Poincar\'e discussing the
doctrine of mathematical probability, which dominates the inanimate world.

\begin{extract}
``You wish me to tell you about these complex phenomena. If by ill luck I
happened to know the laws which govern them I should be helpless. I should be
lost in endless calculations and could never supply you with an answer to your
questions. Fortunately for both of us I am completely ignorant about the
matter. I can therefore supply you with an answer at once. This may seem odd.
But there is something odder still, namely, that my answer will be right.''
\end{extract}

It is perhaps fortunate that we know nothing about the ultimate nature of the
fundamentals of either the physical or mental worlds. We have pursued each so
far as to know that both alike lead away from rather than toward the solution
of the problems of life. The sublimated theoretical concepts in either case
have long ceased to possess actuality. We have rather to find the interaction
between their commonest forms, matter and energy on the one hand and will and
direction on the other.

\origpage{9}

Let us now leave generalities and concentrate upon the question as to what
precisely humdrum mechanical science can contribute to economics. It insists
primarily on the fact that life derives the whole of its physical energy or
power, not from anything self-contained in living matter, and still less from
an external deity, but solely from the inanimate world. It is dependent for
all the necessities of its physical continuance primarily upon the principles
of the steam-engine. The principles and ethics of human law and convention
must not run counter to those of thermodynamics. For men, no different from
any other form of heat engine, the physical problems of life are energy
problems. You have to consider the source, the sunshine. It supplies a
continuous revenue of energy which is consumed by the living engine in its
life. Consumption here does not mean destruction, for destruction, like
creation in the world of which we speak is an absurdity, but merely the
rendering unfit for further use. All the radiant energy received from the sun
sooner or later finds its way into the great energy sink, the ocean of heat
energy of temperature uniform with the surroundings, and is incapable of any
further transformation. This is the form we know most about. It is the energy
of the perpetual thermal agitation of the molecules of which Poincar\'e spoke,
and of which we know nothing (of any individual molecule's motion) and yet
know everything (of the statistics of the motion as a whole). And, it is
useless.

We have next to consider the transformation of the form in which nature
supplies the energy into the form men can utilise and assimilate. In general,
transformation of energy can proceed only in the one direction, much as water
only runs downhill. The water may do useful work on the way turning
water-wheels, or it may not, but may reach the ocean level quite unutilised.
So of the revenue of sunshine which ultimately warms imperceptibly the whole
mass of the globe, it may en route energise a man, or, again, it may not.

As regards the utilisation you have to distinguish very carefully in Cartesian
economics two uses. First, there is the fundamental metabolic use in the body
for the life process, which I shall for brevity term life-use. Secondly, there
is a use in lieu of the first for the doing of external work or labour, better
done directly by inanimate energy. This I shall term the labour-use.

Physically, the life-problem is the reversal of winding a clock. Before any
man can confer the animation of his body

\origpage{10}

\noindent upon a mechanism as in clockwinding, the animation of nature's
mechanism has first to be conferred upon him. History could be rewritten from
the standpoint of how this has been done. At first it was done blindly and
intuitively by trial and error, the survival of the fittest and the
extravagantly wasteful methods that only the unconquerable resurgence of life
can afford. Even now the process is so indirect, being only possible through
the agency of vegetable life, that few realise the terms on which they exist
or the supreme importance of the original sources and amounts of energy
available.

But the labour-use of natural energy has always been a matter of consciously
directed effort and development, since the use of the wind in navigation, and
had proceeded far before the formulation of the principles of energetics. But
in neither case is the sudden break in the continuity of history, which marked
the age of steam, explained by these developments. The key is to be found in
this. Pre-nineteenth century man lived on revenue. Present day man augments
the revenue within certain well-defined limitations out of capital.

All forms of energy previously utilised by life, with one or two minor
exceptions, as tidal energy and that of hot springs, were forms of the solar
revenue. Wind-power, water-power and wood fuel are parts of the year-to-year
revenue of sunshine no less than cereals and other animal foods. But when coal
became king, the sunlight of a hundred million years ago added itself to that
of to-day and by it was built a civilisation such as the world had never seen.

The fundamental fact underlying this civilisation is that whilst men can
lighten their external labours by the aid of fuel-fed machinery, they can only
feed their internal fires with new sunshine and then only through the good
offices of the plant. The vegetable world alone can transform the original
flow of inanimate energy into vital energy. The animal, as yet, is
constitutionally incapable of effecting this transformation.

The technical features of this subject are not without significance. Whatever
the origin of the energy, the penultimate step must always be its storage by
the plant precedent to its use by animals for food. It is possible to draw
upon the energy of a water-fall and to store it up in various chemical
compounds by electro-chemical methods, and in this form to supply it as a
fertiliser to the plant. Increased crops are so produced supporting an
increased population. The reversal of clock-

\origpage{11}

\noindent winding has been consciously achieved. The falling weight of
Niagara's waters work the man. There is no technical objection to utilising
the energy of coal in the same way, other than that of prime cost. But for
practical purposes it is true that the great capital store of energy in fuel
is not yet utilised for the life-use, but only primarily for the labour-use of
energy by life. The life-use demands the intermediary of the plant, and though
coal was once alive it is long since dead. The laborious and wasteful travail,
through farming and agriculture, has once again to be gone through. In spite
of the striking advances of the past century, the agriculturalist, peasant and
farm-labourer form the dominant economic class, and will remain so until some
new discovery of science deposes them. To my mind this is one of the least
obvious and yet most fundamental facts of economics and social science at the
present time.

It certainly has not been sufficiently realised by economists, particularly in
this country. In the flamboyant period of the utilisation of the capital store
of energy in fuel which is now closing, so far at least as this country is
concerned, we could and did by machinofacture make almost every sort of
commodity and all sorts of labour-saving machinery in exchange for the food
which we could not so make and did not make. The population of Great Britain
rose on account of this exchange of capital for revenue, of factory products
for food, from 10\textperiodcentered5 millions in 1801 to
40\textperiodcentered9 millions in 1911. Whereas in Ireland, which has not
coal, it fell from 5 to 4\textperiodcentered3 millions over the same period.
Cartesian economics is capable of diagnosing instantly the root of the Irish
trouble, as Sir Leo Chiozza Money has pointed out.

By this process of exchange of factory products the whole world gradually drew
more and more for its labour-use upon the capital energy of fuel, and used it
to widen the area under cultivation and to transport the harvests from the
most distant regions of the globe and so indirectly augmented the revenue of
sunshine upon which it is still entirely dependent for its life-use.

But this is a very passing phase. New countries grow old. Their populations
tend to expand to the limit of their food supply, and their industries and
manufactures become developed by the aid of their own resources. For a double
reason, therefore, the flamboyant period of prosperity through which

\origpage{12}

\noindent Great Britain has passed is destined to be short-lived.
``Imperialism'' marks its final bid for survival.

Coal is the real capital, out of the consumption of which the capitalist
civilisation has been built up, but, as regards the means of livelihood of the
swollen population that has accompanied its exploitation, its use in this
respect has been indirect and will cease. This is the great paradox of
Capitalism. It is capitalistic as regards the accessories, conveniences and
luxuries of existence. As regards its necessities it is still, to coin a word,
revenual. Even Adam Smith could say, ``When food is provided it is easy to
find the necessary clothing and lodging.'' To-day, by the development of
mechanical power, it is vastly easier than then. But, once this has done all
it can to develop new countries and increase the food supply, it can meet the
demand for bread only by offering a stone. True, by the advances of chemical
and biological science, by the development of agricultural chemistry and the
breeding of better brands of wheat, much may be done, but scarcely as much as
will provide for the requirements of a four or five-fold increase of
population.

The industrialised countries are, with an enthusiasm reminiscent of a lunatic
asylum, turning out an ever-increasing plethora of mere factory products and
sending them forth to compete in ever-shrinking markets in exchange for food,
and are pouring forth an ever-increasing stream of armaments to fight amongst
themselves for markets. The only goal in sight is war and yet war, the blowing
up of the plethora and the permanent devitiation of the stock of the white
race, at the time, too, when, by reason of failing fecundity, the prospect of
its having to fight about something other than markets is becoming evident.

Physical science thus answers precisely, and, I think for the first time, the
problem of political economy, or, as one Marxian writer puts it, ``What are
the sources of our society's wealth, that is, the means of subsistence and
comforts of the individuals comprising it?'' The means of subsistence are
derived from the daily revenue of solar energy, through the operations of
agriculture. The accessories of life, clothes, houses and fuel, as well as its
comforts and luxuries, are derived in great part by the augmentation of this
revenue out of a capital store of energy preserved from bygone geological
times. Life depends from instant to instant on a continuous flow of energy,
and

\origpage{13}

\noindent hence wealth, the enabling requisites of life, partakes of the
character of a flow rather than a store.\footnote{Had Karl Marx lived after
instead of before the establishment of the modern doctrine of energy there can
be little doubt that his acute and erudite mind would easily have grasped its
significance in the social sciences. As it was, in fairness to him it must be
said that he did not attempt to solve the real nature of wealth, but
concentrated entirely upon the problem of its monetary equivalent, that is,
upon exchange-value rather than use-value.\par\medskip
Lest I be misunderstood I may emphasise here that I am using the term energy
throughout these lectures in the strictest scientific sense, for potential or
kinetic energy as understood by the physical scientist and engineer, and never
in the vague and misleading sense of ``mental energy'' which I have rather
termed the guidance and direction of physical energy. In this strict sense of
the word it cannot be maintained that wealth wholly originates in human
labour, for there is no real distinction in physical science between animate
and inanimate energy. But since wealth is not available energy merely, but
rather available energy usefully directed, or some embodiment of it, human
labour (that is, some form of intelligent human activity which may need only a
minimum of physical energy) is usually, though not necessarily, an essential
factor in its creation.}

This answer, though of fundamental importance to social science and to
political philosophy, has little application to present economical systems,
because these are founded upon a simple confusion between wealth and debt, or,
to put it another way, between the wealth of the community and the wealth of
the individual member of the community.

The wealth of the community is its revenue, which, in the last analysis, is a
revenue of energy available for the purposes of life. That being given, in
sufficient amount and in form capable of being utilised by the existing
knowledge of the time, everything requisite for the life of the society can be
maintained. It is impossible to save or store this flow to any appreciable
extent. True, you can dam a river, at great expense, and make a reservoir.
But, even if not used, the accumulated waters evaporate and leak away. You can
under the same, but even more unfavourable terms, store electric energy. But
to contemplate storing wealth on a national scale for even a day is something
like contemplating a storage battery large enough to satisfy the demand of the
world for electric power for one day. True, nature has stored it in coal by
processes requiring geological epochs, but what we do is to unstore it, an
easier matter, and to convert it into a flow before it is of the least
possible use to us. Again, for short periods, the flow may be embodied in

\origpage{14}

\noindent some concrete commodity, in food which rots, in houses which fall
into desuetude if not kept perpetually under repair, and in all the tangible
assets of our civilisation, in railroads, roads, and public works, factories,
wharves, shipping and the like. All alike are subject to a process of compound
decrement, needing ever larger annual expenditure of new wealth to maintain
them in order, and even then rapidly, with each fresh advance of science,
becoming out of date. Such accumulated assets, at best, are classified not as
accumulated wealth, but as aids and accessories in the maintenance and
increase of wealth out of the available revenue of energy. The wealth is the
revenue, and it cannot be saved.

The wealth of an individual, on the other hand, is something totally
different. The ordinary modern individual member of the community in the vast
majority of cases does not possess enough wealth to keep him alive for a week.
By means of a token, legalised as a form of currency, whether a cowrie stone
or a metal counter, but now, more and more exclusively, a simple paper note,
the community acknowledges its indebtedness to the holder of the token, and
empowers the individual to indent upon the revenue of real wealth flowing
through the markets at any time. Even at this stage we see the interests of
the community opposed diametrically to those of the individual members. As
Ruskin puts it, it is the rule and root of all economy that what one person
has, another person cannot have, and the more the private individual is able
to indent upon the revenue the less is left for public services and the
carrying on of enterprises designed to increase the revenue for the general
benefit rather than private profit. The concern of the scientific man is with
the revenue and with how it may be increased in the directions most essential
for the general well-being. If and in so far as political economy can claim to
be a science, that, also, should be its first concern. The individual, on the
other hand, is concerned only to obtain a larger share of the revenue for his
own private use. In so far as he can only do so by increasing, or aiding in
increasing, the real revenue of wealth for the purpose of use rather than of
usury the community gains. A very real complaint that the worker has with the
existing system is that it provides much more easy and lucrative means of
making money without any contribution to the general wealth, and sometimes
actually by destroying it, by individuals possessed of sufficient of the power
conferred by money to hold up the revenue for usury.

\origpage{15}

Ruskin appears to have had a very much clearer conception of the real nature
of wealth than either earlier or later economists. He pointed out, and his
view would now be understood by anyone who has suffered from the dearth of
servants on account of the war, that the art of becoming rich was to get more
relatively than other people, so that those with less may be available as the
servants and employees of those with more. In this acute and original analysis
of the real nature of the individual's wealth---power over the lives and the
labour of others---Ruskin disclosed probably the most important difference
between the interests of the individual and the interest of the State, and the
main reason why the mastery of man over nature has hitherto resulted in so
meagre a contribution to the perfection of human life. For this reason the
community in its struggle with nature resembles an army officered almost
entirely by the enemy. Of what use are the discoveries of scientific men into
new modes and more ample ways of living so long as the laws of human nature
turn all the difficultly won wealth into increased power of the few over the
lives and labours of the many?

In another respect Ruskin was vastly ahead of his own, not to say our own,
time. He and Marx both fully appreciated the main contention of present day
exponents of economics from the point of view of the creator and producer of
wealth rather than that of the financier or merchant. The wealth of a
community can only be increased by production and discovery, not by
acquisition and exchange. In commerce and exchange ``for every plus there is a
precisely equal minus.'' But the pluses wax magnificent and the minuses
``retire into back streets or underground,'' which renders the algebra of the
science peculiar.\footnote{``Unto This Last,'' John Ruskin, 1877.}

This, then, is my main quarrel with orthodox economics, that it confuses the
substance and the shadow. It mistakes debt for wealth and is guilty of the
same mistake as the old lady, who, when remonstrated with for overdrawing her
account, promptly sent her banker a cheque for the amount.

The confusion enters even into the attempt of the earlier economists to define
the main subject matter of their studies---``Wealth,'' though the modern
economist seems to be far too wary a bird to define even that. Thus we find
that wealth consists, let us say, of the enabling requisites of life, or some-

\origpage{16}

\noindent thing equally unequivocal and acceptable, but, if it is to be had in
unlimited abundance, like sunshine or oxygen or water, then it is not any
longer wealth in the economic sense, though without either of these requisites
life would be impossible.

Now it is the object of science to render the enabling requisites of life,
such as food, warmth and other forms or embodiments of energy necessary for a
decent existence, so abundant that they shall cease to be wealth in the sense
of the economist. By increasing a real quantity you do not diminish it, nor by
increasing it without limit do you destroy it. The object of science is to
destroy wealth in the economist's sense of debt altogether by increasing real
wealth without limit.

At the first blush and before they have had time to think, most youthful
students of economics will probably tell me that I am playing with words by
using the word wealth in two senses equally well understood by the economist.
The fact is that the economist, ignorant of the scientific laws of life, has
not arrived at any conception of wealth, apart from the elaborate code of
enactments and legal conventions which give to the individual in actual
non-possession of wealth the right to acquire it, whereas I, from the
application of the laws of energy to the problem of how men live, have arrived
at such a conception.

In conclusion, I may devote my attention to the commonest form of debt, money,
because I believe that until correct views are more widely diffused about this
convention, and the purchasing power of money is fixed as definitely as are
the standards of weights and measures, there can be no peace in society and
the whole elaborate political and social system will remain merely a dreary
and elaborate make-believe.

Once more owing to the war, the real nature of money can be apprehended by
anybody. It ought to bear precisely the same relation to the revenue of wealth
as a food ticket bears to the food supply or a theatre ticket to a theatrical
performance. Whereas, as a matter of fact, at present there is no more
connection between the currency and the revenue than there is between the
birth-rate and the barometer. The revenue depends upon the chances of the
harvest, and all the causes such as the prevalence or absence of disease,
tempests, drought and sunshine which affect the productivity of nature. The
currency is, or was, left to the luck of the gold prospector, and his
spasmodic discoveries, to the state of knowledge of extracting the precious
metals in which a single innovation, such as

\origpage{17}

\noindent cyaniding, may enormously increase the supply, the invention of such
a system as that of cheques, the solemn carting around of gold from one
capital to another to concertina the prices up and down to suit the hierarchy
who have made of money a mystery and of the currency a never-failing
confidence trick.

Whereas, if money is to fulfil its function as a measure of value, it is clear
that the currency must be regulated \textit{pari passu} with the changing
revenue, issued as the latter expands and destroyed as the latter contracts.
Since it would neither be given away in the first nor taken away in the second
event, but used to buy back old, or taken in exchange for new State loans, the
community as a whole would share the prosperity of good times as well as the
stringency of bad ones, instead of only the latter as under the existing
system.

I remember reading as a young man, in some book on economics which I have not
since been able to trace, of the almost mystical virtues of gold in human
welfare and how each successive discovery of that metal in California, South
Africa and Australia was followed by a boom of trade and increased national
prosperity. To a chemist the mystical virtues of any metal, even gold, seemed
a wholly chimerical illusion, but I waited twenty years before the real
explanation became obvious.

Last century was a time, when, wholly beyond the understanding of those who
lived in it, science was increasing the revenue of the world by leaps and
bounds by the consumption of the store of energy preserved in coal. If the
food supply is increased without a corresponding issue of new food tickets,
every holder of an old ticket gets proportionately more. Whereas if the ticket
issue is increased \textit{pari passu} with the food supply, the old ticket
holders get the same as before and fresh people get the surplus of food. Hence
every increase of currency in that flamboyant era of prosperity, whether it
resulted from the discovery of gold-mines or the invention of cheques, meant
that the increased prosperity did not go to the community's creditors, but a
part corresponding with the increase of currency went to fresh people and
general prosperity was the result. How much easier it would have been simply
to print the money and use the issue to repay the National Debt. But the
opportunity passed and the like may not occur again.

I shall be asked by those who are unaware of the proposals made by Gesell on
the Continent and by Kitson in this country how is it possible to fix the
purchasing power of money. The

\origpage{18}

\noindent answer is simple enough. By fixing it, that is, by printing more as
average prices, determined by index numbers, tend to fall and by withdrawing
it from circulation as they tend to rise. As it is, these matters, which are
the most vital factors of all that enter into the economic welfare of the
community, are left to the oddest combination of natural luck and human
cunning to which, surely, any race ever entrusted its destinies.

Money, I shall be told, must function not merely as a measure of value, but as
a medium of exchange and as a store of value.

As regards the latter, humanity is crying for the moon. Wealth is a flow, not
a store. After the search-light of the war, I can conceive no nation so
barbaric as to regard gold as a store of value. Demonetise it and where is its
value? Not a gold mine would be at work on the morrow. The world has enough
gold to stop its teeth and gild the inside of its tea-spoons for hundreds of
years. Nor, as a medium of exchange, can anyone, after the experience of the
war, really find any fault with paper, provided of course its issue were
directed to the end of maintaining average prices constant from century to
century.

Civilised nations maintain at great expense elaborate testing institutions to
fix with meticulous accuracy and disseminate replicas of all quantities that
enter into one side of every commercial transaction involving buying and
selling. They maintain an army of officials and inspectors to suppress the
hollow pound weight, the elastic yard-wand and the telescopic quart pot. What
an elaborate fraud upon the public all this one-sided passion for accuracy is!
The public are not interested in the absolute magnitude of weights and
measures. What is solely of practical importance is the relative measure, not
merely how much coal there is in the sack or how much beer in the mug, but how
much coal and how much beer for how much money.

Do we keep a National Economic Bureau to stabilise the purchasing power of
money and an elaborate organisation of inspectors, the counterpart of those
who suppress petty fraud, to deal with organisations to concertina the pound
sterling? Our system is precisely analogous to sealing up only one arm of a
balance and making an imposing parade of protecting it from the wind and
tamperers, while leaving the calibrating arrangements of the other arm to the
manipulation of a class of persons deriving their livelihood from the
business. It is

\origpage{19}

\noindent on record that a group of American financiers on one occasion,
having sold British and bought American securities in advance, removed
\pounds11,000,000 from the Bank of England and put it into circulation in
America, with the result that the prices of the securities they had sold fell
greatly in value and those they had bought rose correspondingly. Since gold
never constitutes more than a few per cent.\ of the total currency, a
reduction or increase of the gold basis in a country is followed by an
enormously larger total fall or rise of values, and financiers in the position
to cart about a few millions of the ``precious'' metal at their will can very
easily and certainly acquire other people's wealth.

No doubt the instance quoted is an extreme one, but when one enquires further
as to who is in charge of the calibration arrangements that fix the purchasing
power of money, much that has hitherto seemed inexplicable about our time
becomes clear. These powers are wielded, by private banks, like the Bank of
England, in the steadfast interests not of the community but of the creditors
of the community. Whereas no changes of revenue, so long as the currency
remains constant, affect the relative proportion of the whole revenue secured
by the creditors, any increase of currency diminishes their relative share and
hence is known as inflation, while any decrease increases their relative
share, and hence is called sound finance.

Even so far as we have yet got in disentangling current misconceptions of
wealth from reality, it is not difficult to understand why the blessings
conferred by science have been of so limited incidence. Civilisation has been,
in its most vital interests, not in the hands of those who have contributed
most to its wealth, but of those to whom in a very literal sense it is
indebted, and is likely, under this system, to become ever more indebted. This
gives a short practical remedy for the most obvious of the ills that
civilisation is heir to. Institute a complete organisation for ascertaining,
on every public proposal, the feeling of the City, and the views of the
captains of finance and banking, but act in precisely the opposite direction.
From the point of view of the welfare of the community rather than of its
creditors, you could hardly fail to be right every time.

\origpage{20}

\vspace*{4pt}
\begin{center}
  {\addfontfeature{LetterSpace=6}\footnotesize SECOND LECTURE.\,---\,Chairman, Principal Senter.}
\end{center}
\vspace*{6pt}

\lettrine[lines=2,findent=2pt,nindent=0pt,slope=0pt]{S}{OME} questions I was
asked at the end of last lecture seem to indicate the necessity of first
clearing away some misconceptions, partly, perhaps, due to my citing Ruskin as
an economist. Although my views are very similar in some respects to those
arrived at long ago by Ruskin, I may be permitted to remark that I have
deduced them, without at the time being aware of Ruskin's writings on this
subject, from the principles of the heat-engine, rather than from those of
ethics. I know it is a burning question whether economics ought to concern
itself with ethics at all, but of its obligation to understand the engineering
of life I do not think there can be two minds. If it is a science at all, it
is, in Huxley's words, concerned with truth as ``veracity of thought and
action, and the resolute facing of the world as it is when the garment of
make-believe with which pious hands have hidden its uglier features has been
stripped off.'' Neither the ethical nor statistical sides of make-believe are
to-day of any very great interest, but economics has still to achieve the
emancipation which, in Huxley's day, the biological sciences accomplished. It
is just because the application of the every-day principles of engineering to
the living engine offers such a powerful corrective to the make-believes of
the economic systems of society that I have ventured to address you on the
subject.

On the strength of a quotation from Ruskin, that there was no profit in
exchange, Ruskin was condemned by one distinguished economist, I think
unfairly. I am well aware that it is the fashion to regard Ruskin as
out-of-date as an economist, though as a matter of fact the times are even yet
hardly ripe for a judgment on this point. But on the actual statement that
there can be no profit in exchange there can surely be merely a difference of
meaning to be attached to the word profit.

There is much making of money but no making of wealth by exchange, much
advantage to Society for which the merchant is highly remunerated, much
acquisition by the merchant of wealth, but no profit, for the sum total of
wealth is unaffected by exchange and against the merchant's plus there is a
precisely equal minus.

The word economics was coined by Aristotle as signifying household management
in contradistinction to money and trade (chrematistics). What Aristotle meant
2,250 years ago,

\origpage{21}

\noindent I pointed out again last lecture when I charged ``economists'' with
confounding debt for wealth. A ham merchant working on what he is pleased to
call a 10 per cent.\ basis of profit, may buy ten hams for the same sum as he
sells nine. He may be pleased to think he has made a profit of one ham, but he
certainly has not made a ham. There were and remain ten, whereas if anyone had
made a profit of one ham, there should now be eleven. These hams represent the
life-time profit of a certain number---2 to be precise---of pigs, fed,
according to nursery tradition, on the skins of potatoes, which in turn
derived their feeding value from the sunshine. Wealth being some form of
embodied useful energy, the law of the conservation of energy applies to
wealth in that for every plus there is a minus. But fortunately in this case
the earth is credited with the plus while the sun is debited with the minus,
and that is as good as an actual creation of wealth from the terrestrial point
of view. Nearer than that the laws of matter and energy do not allow.

The opposite (for every minus there is a plus) is not true of wealth, because
we are dealing with the availability of energy rather than with its total
amount, and because of the natural tendency of all available or wealth-forming
energy to pass, more or less quickly, into the waste heat of uniform
temperature of the surroundings.

The exposure of the shams of economic systems is a time-consuming task, and,
as we have still to consider the nature of capital and usury, it is as well
first to say something about the realities. How is wealth produced and what,
if any, are the limitations to the wealth of an intelligently directed
community?

The first factor, a continuous flow of energy of an available form, has so far
only been considered. If that were unlimited in amount and under human control
in the same way as the energy of fuel now is, this factor would impose no
limit on the production of wealth. Even the distinction which it is at present
so necessary to make, between the life-use and the labour-use of energy, would
be of less importance, for it is not so much that the synthetic production of
food-stuffs, except by the aid of the plant, is impossible as that it is
impracticable with energy at its present value. Last century it must have
appeared, to any one following the line of thought we are pursuing, that the
limitations of this first factor of wealth must

\origpage{22}

\noindent always limit human ambition and expansion. But now we know that it
is not so. The extraordinary developments since the beginning of the century
in the study of radioactivity and of the internal structure of the atom have
proved that there is resident in ordinary materials amounts of energy of the
order of a million times that which can be obtained from fuel during
combustion, but that to liberate this store the transmutation of the elements
one into another must first be made possible. The radioactive elements are in
course of a natural transmutation, which, while it is impossible to stop, is
likewise impossible to imitate. The energy of radium, the element which for
thousands of years emits as much heat every two days as its own weight of fuel
in burning, is derived from this hitherto unsuspected store of energy in the
structure of the radium atom in its change into atoms of lead and helium.

No! it is the second factor in the production of wealth that now limits, and
probably will always limit, human prosperity. It is knowledge, or rather
ignorance. For untold years men froze on the site of what now are coal mines,
and starved within sound of the Niagara that is now at work providing food.
Every single factor in wealth production existed prior to the phenomenal
expansion of the last century except one, the knowledge how to control and
utilise for life the capital store of sunlight preserved in fuel. It is
precisely the same to-day. We are as far from utilising the stores of energy,
which we know exist all round us in unlimited abundance, as savage men, who
had not yet learned how to kindle a fire, were from utilising the power which
has made our own age great. The whole matter could not compete in public
interest with a ball game or a prize-fight, and, as has been recently said,
civilisation depends for its future on the long vacations when the scientific
men in the Universities get the opportunity for a few weeks' uninterrupted and
continuous research.

Although it is far from my own view of the matter to divide, as is sometimes
done, the winning of knowledge into two water-tight compartments, pure and
applied, academic and technical, or discovery and invention, and to elevate
the former upon a pinnacle attainable only by the few and to depress the
latter to a level only a little above the capacity of the ordinary efficient
routine worker, it is undeniable that in point of time at least, pure
scientific knowledge, acquired for the sake of knowledge only and with no
definite utilitarian end

\origpage{23}

\noindent in view, must invariably precede any great advance in technology and
invention. But in both fields the qualities of mind and temperament required
are much alike, and in both, in their highest expressions, attain that
undefinable and elusive quality we characterise as genius. All genius in this
respect is alike---it creates, and in all creation the whole is invariably
incomparably greater than the sum of the component parts.

Those whose philosophy consists in passing by little steps, each almost
trivial apparently, from the electron to the soul, try also to pass from the
humblest beginnings of intelligence above the animal level to the present
heights of intellectual achievement attained in the exact sciences. Genius to
such is ``only'' the cumulative sum of infinitely little steps in intellectual
progress which began with man himself. To my mind you might as well describe
an old master as a cumulative effect of infinitely little daubs of paint, or a
symphony of an accumulation of sound vibrations. The whole is greater than the
parts, but even so, that is not the chief point that is missed. Everyone knows
the difference between reading or translating a foreign language and speaking
it. Some of us, who are being required now to spend a year of research work
before taking the degree may realise the difference between knowing all about
every important advance ever made in the subject of our study as well as or
better than those who actually made these discoveries, and achieving the most
infinitesimal advance therein ourselves.

Just as I am constrained to put a barrier between life and mechanism in the
sense that there is no continuous chain of evolution from the atom to life, so
I put a barrier between the assimilation and the creation of knowledge. Each
one of these infinitesimal steps of intellectual progress which look so small
in retrospect, once had a different character. Otherwise why, for example, was
it left for Newton to discover the law of gravitation or Benjamin Franklin the
nature of lightning?

We have here a rather remarkable peculiarity of the human mind. Every teacher
knows how apt his best pupils are to acquire and follow the achievements of
the past and present, so that at the age of twenty they often may have a
vaster range of knowledge than any of the pioneers who contributed to the
subject. Yet how rare, so far at least, is it to find these uniquely equipped
and finished repositories of knowledge capable of making a single step
forward, that is not merely imitative, and

\origpage{24}

\noindent which a hundred years hence will survive to the honour of appearing
trivial? Possibly we are on the eve of understanding something about the
nature of intellectual creation as distinct from mimicry. But until we do, the
wealth of the world, material no less than spiritual, has, ultimately, to be
ascribed to the work of singularly few minds.

After energy and genius we come to Labour, the one factor in wealth production
which hitherto has been adequately recognised. So far at least this factor has
not limited the expansion of wealth, but rather expands \textit{pari passu}
with it, though, in this respect, in all the European countries and also
Australia, the increase of population appears definitely to have been checked.
It is probable, at least of our own day, that any definite retrogression, as
regards knowledge once attained, is unlikely, and even if civilisation has to
pass through a time in the future analogous to the Dark Ages of the past, the
intellectual achievements of the present day will remain preserved. Granted a
certain stage in the technology of wealth production, it is probable that
this, without any further aid from what I regard as the essentially creative
type of mind, can be maintained from generation to generation by Labour in its
widest sense, including in that term every kind and grade of routine and
imitative worker, whether by hand or brain.

Thus we have three factors, two of the character of a continuous never-ending
contribution, a flow of energy and the unremitting attention, whether physical
or mental, required for its utilisation, and one, the creation of knowledge,
of the character of a definite step forward made once and ever afterwards
available for all time. The latter factor limits the rate of expansion of
wealth but, strictly speaking, contributes nothing to its actual production.
It will ever be the prerogative of genius to endow posterity rather than its
own day.

More and more as time goes on, this character of wealth, essentially as a flow
rather than something that can be stored, forces itself upon our attention as
we grow out of the merely subjective views we derive from our own means of
livelihood and bank balance. The time is fully ripe that the world should
reconsider from a scientific basis the conventions by which it empowers
individuals to ``save'' and amass ``riches.'' Broadly, the object of such
conventions should be definitely directed to provide for the period of
childhood, adolescence and old age, for genius and for the dissemination and
diffusion

\origpage{25}

\noindent of its results and for similar work of a publicly beneficial
character. These must be charges on the revenue rather than debts handed on by
private individuals to their heirs and successors and swollen by the
ridiculous pretensions of the usurer to an absurdity.

With the growing appreciation of the general public, made wise because of the
war, of the real nature of money, the much vexed question of what is usually
called capital should not give us much difficulty. Just as money is a paper
indent upon the revenue, capital is the paper receipt for the expenditure of
wealth. Economists brought up on the mythical origin of man, as recorded in
the book of Genesis, used to be fond of inventing, to explain the origin of
capital, a mythical Robinson Crusoe, of exceptional industry and acumen, as
the primitive capitalist. With the advance of knowledge the real Adam has
turned out to be an animal, and now the original capitalist proves to have
been a plant!

The material and scientific greatness of our day is due to the primitive
accumulation of the solar energy of the forests of the carboniferous era, and
preserved to this day as coal. The plant accumulated, we spend.

When coal is burnt it is burnt. You cannot both burn it and keep it in the
cellar, and still less can you go on drawing interest from it for ever at so
much per cent., as is the case with the so-called capital of the economist and
the business world. Here again the economist is mistaking our old friend debt
for wealth. The wealth has been spent, not saved, and exchanged for some form
of receipt, giving the holder a purely conventional right to so much per cent.
per annum until the debt is repaid.

The capitalist wishes to have it both ways, to be regarded as a public
benefactor because he spends his wealth, not in drinking himself to death, but
in enterprises designed to increase the revenue. If he did this he would
indeed be a public benefactor. But the community having spent his wealth, as
regards himself he expects it all back in due course with interest on the
loan. The consequences of his abstinence are that civilisation has got
inextricably ``into the hands of the Jews.'' Compared with this, the wildest
profligacy on the part of the original capitalist would have been a relatively
minor evil.

It is of course a colossally hard task wisely to expend

\origpage{26}

\noindent wealth so as to increase the revenue. All honour is due to the
business energy and enterprise of the commercial and technical managers and
the workers who effect the conversion of capital wealth into increased
revenue. But as regards the people who merely lend the wealth spent, the
ordinary dividend holder of a joint-stock company for example, he is of course
simply that peculiar type of benefactor which used to be termed a usurer. We
are all in it now, ever since it became possible to buy a \pounds1 War Saving
Certificate bearing compound interest for 15s.\ 6d. The extraordinary changes
of legal and social conventions with respect to interest and usury, recorded
in history, make it quite clear that political economy, which depends upon
such factors quite as much as upon the laws of energy, can never be a science
in the same exact sense as physics or chemistry. To Aristotle a usurer was a
person beneath contempt. To-day, even the Vice-Chancellors of the ancient
Universities, which purport to hold up to reverence Greek thought and culture,
are as enamoured as anyone of the excellence of compound interest.

Of all hard critics of the usurer, Martin Luther is easily first, and in his
vigorous denunciation there is a certain perspicuity which we moderns seem to
lack. Otherwise few could tolerate the economics of an ordinary daily
newspaper or social club.

\begin{extract}
``The heathen were able by the light of reason to conclude that a usurer is a
double-dyed thief and murderer. We Christians, however, hold them in such
honour that we fairly worship them for the sake of their money.

``Usury is a great huge monster, like a were-wolf, who lays waste all, more
than any Cacus.\ .\ .\ .\ For Cacus means the villain that is a pious usurer
and steals and robs and eats everything. And will not own that he has done it
and thinks no one will find him out, because the oxen drawn backwards into his
den make it seem from their footsteps that they have been let out. So the
usurer would deceive the world as though he were of use and gave the world
oxen while he however rends and eats all alone.''
\end{extract}

One must admit that it would be difficult to find a better description of
usury than is given here, ``Oxen drawn back into the den which, from their
footsteps appear to have been let out.'' Current orthodox economics gives the
credit that

\origpage{27}

\noindent rightly belongs to the scientific discoverer, the expander of
wealth, to the usurer, the expander of debt.

I shall be told that there is something to show for capital expenditure and
that against the paper receipts there are tangible assets. Thus, if a railway
is taken as an example, there are the rolling stock and the rails. But
admittedly these would have but a scrap metal value if railways ceased to pay
and the shareholders wanted their money back. As the world gets older all this
initial expenditure has to be periodically renewed to make good depreciation,
and, further, the plant and the methods of operation become, with the advance
of knowledge, out of date. The indebtedness to the original shareholders does
not, however, usually cease on that account. Railways continue to pay
dividends on all capital expended, though, as in the case of the canal systems
purchased, much of it altogether ceases to bring in revenue. It seems to me to
be merely a matter of time before this happens with every form of capital
expenditure. The normal old-age form of capital is simple debt, a permanent
lien upon the future revenue of wealth. The assets are much overrated. If the
world worked as hard on construction as it did during the war on destruction,
and was permitted by the usurer to do so, all civilisation could probably be
rebuilt upon an up-to-date plan and the Augean stable of a modern
industrialised community cleaned up in less time than the war took.

The vast heritage of wealth which science made available at the commencement
of the 19th century, in so far as it has been spent, has been replaced by
paper receipts for the expenditure which bear interest in perpetuity. Capital
merely means unearned income divided by the rate of interest and multiplied by
100. If I invent a process bringing in \pounds1,000 a year revenue, its
capital value is \pounds20,000 if the rate of interest is 5 per cent., and I
can sell it for some such sum. The capital of the world in this sense to-day
aggregates to an altogether inconceivable sum. There never has existed at one
time such an amount of wealth. It represents the accumulated capital
expenditure of generations of men. During the war the capital of the country
was increased by some \pounds7,000,000,000, which brings in \pounds350,000,000
a year permanent interest. I may be told that everyone will admit that this is
debt, whereas it is in fact, in this respect, precisely on the same footing as
so-called productive capital expenditure, a private lien upon

\origpage{28}

\noindent the revenue, wealth to the individual owner and debt to the
community. As regards the national accounts this \pounds350,000,000 a year is
a simple transference, rather than an item of expenditure. The sum is
collected from the taxpayer, and paid to the holders of War securities. It
amounts to \pounds8 13s.\ 4d.\ per head of population, and includes the small
item of 10s.\ per head in respect of the Napoleonic wars. There can be only
one possible end to this process. Though for a time the advances of science
may so increase the revenue from year to year as to render these payments by
way of interest possible, in the end the whole of the revenue must be in the
control of the usurer. A small part of the population will get into the
position of a great rentier class living on interest, and most of the rest
will be reduced to starvation in so far as they are not kept alive by State
doles. How far this process has already gone in this country is obvious, since
something like a quarter of the population at the present time is unemployed
and the expenditure on national education is only about a quarter as much as
that upon the holders of war securities.

Wealth is a flow and it cannot be saved. Spent it must be as it accrues,
whether on consumption or on capital outlay designed to produce future wealth.
As regards the first, life is consumption from the cradle to the grave,
consumption of that pristine flow of energy we owe to the sun. The efforts of
the financier and monied person to make life a balance-sheet with the debit
and credit sides in agreement are untrue to nature. Life is a continuous
expenditure of wealth and in this point again Ruskin, rather than the modern
chrematist, was absolutely in agreement with science. And, as regards the
second, capital expenditure, even though it is designed to increase the flow
of wealth, and admittedly, over a certain natural period which is not
infinite, it does achieve this object, it is expenditure as much as the other.
The river of wealth is thus divided, but a part is metered carefully and
recorded as an accumulation of capital indebtedness against the community.
Scientific men, in the innocence of their hearts and the benevolence of their
souls, fondly imagine that by increasing the revenue available for life they
benefit the community. But do they? The larger the revenue, and the greater
the flow above that required for immediate consumption, the greater the debts
incurred by the community and the impossibility of their ever being repaid.
Meantime, though real

\origpage{29}

\noindent wealth rots if stored, the meter readings spontaneously bear
interest and increase \textit{ad infinitum}. The principles and ethics of
human conduct and convention have their own code and standards, but whatever
they may be, they must conform to and not run counter to the principles of
thermodynamics. A chauffeur may have a soul above the mechanism of his car,
but if it led him to try to run it on already consumed petrol, none the less
he would be considered a great ass.

The usurer in Martin Luther's time had not achieved the colossal success in
``deceiving the world as though he were of use and gave to the world oxen''
that he can claim to-day. The professional economist seems to be an easy
conquest. Thus Mr.\ J. M. Keynes, in his ``Economic Consequences of the
Peace,'' seriously seems to think that the law of compound interest is the law
of increment of wealth rather than that of debt, and offsets it against the
Malthusian law of increase of population! ``One geometrical ratio might cancel
another and the 19th century was able to forget the fertility of the species
in a contemplation of the dizzy virtues of compound interest.'' To him capital
is a vast accumulation of fixed wealth, in danger of being prematurely
consumed in war. He likens it to a cake, which one day, owing to the dizzy
virtues of usury, may be large enough to go round. ``In that day overwork,
overcrowding and underfeeding would come to an end, and men secure of the
comforts and necessities of the body could proceed to the nobler exercise of
their faculties.'' Cake happens to be the one material of which it has been
well said that you cannot eat it and have it too, and I would suggest that
this is the real reason for Mr.\ Keynes's somewhat mystical references to a
peculiarity of capital considered as accumulated cake, that this ``is only in
theory---the virtue of the cake was that it was never to be consumed.'' In a
similar vein we have Major Douglas and Mr.\ Orage advocating the salvation of
the economic system by the bringing of the dizzy virtues of dividends within
the reach of the many rather than of the few. I do not mean to imply that
their system might not be a great improvement upon the present one. Indeed,
any system almost must be an improvement upon one that is administered solely
in the interests of the creditors, rather than the creators and consumers of
the community's wealth. But my analysis would lead me to class their
suggestion rather as a temporary palliative, for no system founded upon usury
can be stable.

\origpage{30}

It is wonderful how people, who never come up against reality from the cradle
to the grave, and live all their lives a purely artificial existence in some
city divorced from all contact with primitive nature, get into the habit of
supposing that the conventions which regulate their businesses and livelihoods
can be applied to the economy of the world at large. It would be quite
impossible for any member of the agricultural community, for example,
accustomed to the ways in which wealth is really produced, to fall down and
worship the institution of usury in such a naive fashion, or for a Labour
Government, to be guilty of the confusions between wealth and debt which are
characteristic of orthodox politicians at the present time. We have in the
hitherto association of the function of Government almost entirely with those
who live by rent, interest and profit, and thus take from rather than
contribute to the revenue of the wealth of the community, a further
justification for the view already expressed, that the community in its
struggle for existence resembles an army officered almost entirely by the
enemy.

You cannot permanently pit an absurd human convention, such as the spontaneous
increment of debt, against the natural law of the spontaneous decrement of
wealth. This applies whether it is simple or compound interest which regulates
the increment of debt. For obvious reasons the law of compound interest over
great periods of time rarely as yet has fully operated. But the significant
and distressing fact is that this absurd law, with the concentration of money
in the hands of trusts and combines of financiers, now tends to operate more
and more fully every day. Some of you may have heard of the story of the
reward asked of the Emperor of China by the man who taught him chess. It
looked modest enough. He wanted one grain of corn for the first square of the
chess-board, two for the second, four for the third, eight for the fourth, and
so on in a geometrical progression to the 64th square. The story goes that the
first half of the board was easily accounted for, but before three-quarters
had been so dealt with, the Emperor had to cry off, as his couriers came back
to him reporting that there was not such a quantity of corn (actually 23
million tons) in the Empire. If the matter had been pushed to the bitter end,
at the 64th square, the number of grains would have been one less than
2\textsuperscript{64}---just about one million million tons---more than the
present popu-

\origpage{31}

\noindent lation of the world could consume in a period of time longer than
that covered by the records of history. This is the law of compound interest.
To-day \pounds1 of debt doubles itself in about 12\textperiodcentered5 years,
and becomes \pounds1,024 in 125 years and over a million in 250 years at
5\textperiodcentered5 per cent.\ compound interest. If this is to be the
inevitable consequence of scientific men increasing the wealth of the world,
enthusiasm in well-doing may well dry up. Nevertheless, on no other terms than
a perpetual increase of the revenue by scientific discovery can such a system
of ``economics'' be maintained, and, even so, if were the 20th century as
prolific in discovery as the 19th, there is no escape for it from ruin under
the law of usury and the rule of the usurer.

Let us, in conclusion, put in the light of the analysis of wealth and wealth
production which has been attempted, the aspect of this mad system which is
now uppermost in the minds of many thoughtful people, its inevitable end in
world-war. Remember that it is impossible to save. The revenue of wealth must
be spent as it accrues either in consumption or in capital expenditure. The
latter is for the purpose of securing a lien on the future revenue, by
producing more goods or services out of which a profit can be made. But as the
masses do not receive the profit, they cannot buy this increased production.
It was the discovery of the classical economists that wages are the purchasing
power necessary to maintain a supply of labour, that is to provide food and
shelter for the labourer and his family in the particular condition of life
necessary and customary for that kind of labour. Man-power is now cheapened by
being pitted against the infinitely greater and more docile power of inanimate
nature. For a brief period, which is now closing, though only after terrible
hardship and suffering, the displaced labour found an outlet ultimately by
reason of the increased fraction of the revenue devoted to capital
expenditure. But the world fills up. Its markets, at first open to the excess
production of the industrialised nations in exchange for food, tend to close
as time goes on. The fiercest international rivalry for markets ensues, and to
industrialised nations armaments are, as products of machino-facture, the one
thing that can be turned out in almost limitless abundance. But armaments and
war do not produce food. They merely determine the distribution as between
competing nations, and tend to destroy food-producing power to an extent that
makes

\origpage{32}

\noindent even the victors actual losers. In this Ruskin was again far ahead
of his own times. He alone seems to have had sufficient veracity of thought
and power of penetrating below the conventions of society to realise that, in
the frenzied pursuit of gain, the objective was illusory in a physical sense.

\begin{extract}
``Capital is a root which does not enter into the vital function till it
produces fruit. Capital producing nothing but capital is root producing root,
bulb issuing in bulb, never in tulip. The Political Economy of Europe has
hitherto devoted itself to the multiplication of bulbs. It never saw nor
conceived such a thing as a tulip. Nay! boiled bulbs they might have been,
glass bulbs---Prince Rupert's drops consummated in powder---well if it were
glass powder and not gunpowder.''
\end{extract}

I do not pretend to be able to have got further in my own conclusions than
this, that the rule of the usurer in political and social affairs has become
impossible and that he has to go. During a time of expanding revenue, or
before the burden of interest charges on the revenue from accumulated debts
equals the annual expansion, he may be an efficient, if brutal, task-master,
and the lure of private interest and gain may be a safe principle in place of
government. But at a time like the present, when the usurers of the world, if
cheated of their expectations, like Shylock are bent on a pound of flesh next
to the heart, their pretentiousness and futility is obvious, and some form of
government according to economic, rather than chrematistic, principles will
have to be resumed. The laws of energy under which men live furnish an
intellectual foundation for sociology and economics, and make crystal clear
some of the chief causes of failure not only of our own but, I think also, of
every preceding great civilisation. They do not give the whole truth, but, in
so far as they are correct to physics and chemistry, they cannot possibly be
false. I think with but little amplification and modification they might
furnish a common scientific starting point from which all men concerned with
the public rather than with their own private interests might start to rebuild
the world more in conformity with the great intellectual achievements which
have distinguished the present age. The first step towards such a scientific
Utopia would be the due delimitation of the rights of the community's
creditors---the curbing of the demon of debt which masquerades among the
ignorant as wealth.

% ================= COLOPHON =================
\clearpage
\pagestyle{empty}
\null\vfill
\begin{center}
{\addfontfeature{LetterSpace=10}\footnotesize NOTE ON THIS SETTING}
\end{center}
\vspace{10pt}
{\footnotesize\linespread{1.05}\selectfont
\noindent This is a fresh setting of \textit{Cartesian Economics}, first
published by Hendersons of 66 Charing Cross Road in May 1922 and reprinted in
the October of that year. The copy behind it is the Harvard College Library
(Widener) exemplar, barcode 32044010374809, from the bequest of George
Hayward, M.D., accessioned 16 July 1925; the text was recovered from the
optical transcription of that scan and collated line by line against it.\par
\medskip
The folios follow the 1922 printing exactly: what stands on page 12 here stood
on page 12 there, so a citation taken from these pages will answer to the
original. Obvious transcription damage has been mended silently---broken
word-endings at the scanned margin, hyphens carried over from the compositor's
line-breaks, and inverted commas. Two readings remain doubtful and are left as
they were found: \textit{devitiation} (p.\ 12) and the figure of \textit{23
million tons} (p.\ 30). Decimal quantities are set with the mid-point, as a
British house of the period would have set them. The two notes on page 13 are
gathered under a single mark, as the page gives no second reference.\par
\medskip
The text is set in MFB Oldstyle, a digitisation of Century Oldstyle, cut by
Morris Fuller Benton for the American Type Founders Company in 1909 and in
common use for book work on both sides of the Atlantic when these lectures
were given.\par
\medskip
Soddy's language is that of 1921 and is reproduced without alteration,
including the passage on page 25 and the reference to racial stock on page 12,
which the modern reader will want to weigh as evidence of the period rather
than pass over.}
\vfill
\clearpage

\end{document}
